\conclusion

В ходе работы были проанализированы ограничения классических методов регистрации облаков точек в условиях ухудшения наблюдаемости сцены, описан используемый набор данных Boreas-RT и сформирован воспроизводимый сценарий вычислительного эксперимента с моделируемым обгоном грузовика. Разработан алгоритм обработки последовательности лидарных кадров, включающий двойную вокселизацию облаков точек, регистрацию соседних облаков методом ICP, оценку преобразования между парами точек с учётом весов, задаваемых функцией Гемана---Макклюра, а также модель постоянной скорости и фильтр Калмана для стабилизации начального приближения на следующем шаге регистрации. Кроме того, реализован экспериментальный конвейер, позволяющий сравнивать различные конфигурации алгоритма по набору количественных метрик и средств визуального анализа.

Проведённое экспериментальное исследование на 10 тестовых проездах Boreas-RT, нарезанных на 725 десятисекундных сцен, показало, что использование фильтра Калмана не приводит к вырождению оценок скоростей на длинных последовательностях и не подавляет естественную динамику движения автомобиля. В то же время на коротких сценах с моделируемым перекрытием поля зрения применение фильтра Калмана повышает устойчивость алгоритма и уменьшает влияние редких, но сильных деградаций оценки движения. Наиболее содержательный выигрыш наблюдается для продольной компоненты линейной скорости, которая оказывается наиболее чувствительной к ухудшению наблюдаемости сцены.

Анализ результатов показал, что основной эффект от использования фильтра Калмана связан не с формальным сглаживанием выходных оценок, а со стабилизацией начального приближения для следующего запуска ICP. Благодаря этому уменьшается риск последовательной деградации регистрации на соседних кадрах и повышается устойчивость всей итеративной схемы оценки движения. Тем самым предложенный подход позволяет надёжнее проходить короткие критические участки траектории до восстановления большей наблюдаемости сцены.

Вместе с тем проведённое исследование выявило и ограничения предложенного алгоритма. Фильтр Калмана не устраняет фундаментальные трудности, связанные с геометрической вырожденностью отдельных дорожных сцен. В частности, в туннелях и на трассовых участках продольная компонента скорости может оставаться плохо наблюдаемой даже при использовании устойчивой схемы регистрации. Дополнительным ограничением является то, что параметры ковариаций процесса и измерения в данной работе выбирались по инженерным соображениям и не оптимизировались автоматически по экспериментальным данным.

Дальнейшее развитие работы может быть связано с автоматическим подбором параметров фильтра Калмана, уточнением схемы оценки линейной части преобразования между соседними облаками точек, а также с расширением экспериментального исследования на более широкий набор сценариев динамического перекрытия поля зрения. Перспективным направлением также является объединение предложенного подхода с более сложными схемами лидарной локализации, использующими локальную карту или дополнительные источники априорной информации о движении.
